% !TeX program = platex
\documentclass[a4paper,11pt,dvipdfmx]{jsarticle}
\usepackage[margin=24mm]{geometry}
\usepackage{graphicx}
\usepackage{tabularx,booktabs,array}
\usepackage[hidelinks]{hyperref}
\usepackage{pxjahyper}
\newcolumntype{Y}{>{\raggedright\arraybackslash}X}
\renewcommand{\arraystretch}{1.25}
\title{デジタルコンテンツ企画書\\\large 光と影の猟犬}
\author{チーム26\\木澤・村松・中野・森井}
\date{2026年10月16日}
\begin{document}
\maketitle

\section{作品の目的}
\subsection{対象ユーザー}
キャラクターを育てながら、位置取りや役割分担を考える戦略的なバトルを楽しみたいユーザーを対象とする。
素早い操作に依存せず、自分のペースで盤面を確認して判断できるターン制を採用する。
一人で二匹を操作する、ブラウザ向けのゲームとして制作する。

\subsection{提供する体験と設計の考え方}
光と影で性質の異なる二匹の猟犬を育て、戦場の明るさを変えることで、自分の判断から有利な状況を作る体験を提供する。
育成による相棒への愛着と、二匹の連携で課題を突破する達成感を目指す。

この目的に対し、次の三つの仕組みを用いる。
第一に、明るさによって二匹の移動や攻撃の有利条件を変え、光源や遮蔽物を操作する意味を作る。
第二に、敵の次の行動を予告し、その情報から移動・明暗操作・行動順を考えられるようにする。
第三に、クエストの報酬を訓練や装備に使い、育てた二匹で次の課題へ挑む循環を設ける。
これらの仕組みによって狙った体験が伝わるかは、制作後半の試遊で確認する。

\section{作品構想}
本作品は、光の猟犬と影の猟犬を個別に操作・育成する、マス制の戦術ゲームである。
見た目は斜め上から立体模型を眺めるような3Dとし、移動や攻撃の判定は2Dのマップ上で行う。
各ラウンドの開始時に先に動く猟犬を選び、二匹が一回ずつ「移動＋技」を行った後、敵が予告した行動を実行する。

各マスの明るさは五段階とする。
光の猟犬は灯火の移動や光の技、影の猟犬は遮蔽物の設置や影の技によって、盤面の明暗を変える。
光の猟犬は明るい場所、影の猟犬は暗い場所で有利になる。
特に影の猟犬は最も強い光のマスへ入れないため、単に全体を照らすだけでは二匹を有利にできない。

例えば、暗い場所で防御が高くなる影鎧兵に対し、光の猟犬で灯火を動かして防御を下げ、次に影の猟犬で攻撃する。
ただし、敵のマスを最も強い光にすると影の通常攻撃が使えないため、敵と味方の両方に適した明るさを考える必要がある。
また、射撃を予告した灯射兵は、発射元を暗くするか、射線に遮蔽物を置くことで射撃を止められる。
このように、移動先だけでなく、二匹が行動できる条件そのものを作り替える判断を中心に据える。

\section{作品の独創的な点}
\subsection{既存作品との比較}
光と影の操作やキャラクター育成は既存作品にも存在する。
本作品では、それらを「二匹の行動順」「マスごとの明るさ」「固定された敵の予告」に結び付け、盤面を読んで連携する戦術として構成する。

\begin{description}
\item[Shadow Puppeteer]
少年とその影を操作し、光を利用してパズルを解く作品で、一人プレイと協力プレイに対応する\cite{puppeteer}。
本作品では一人が二匹を順番に操作し、光の操作を移動・攻撃条件の変化として扱う。
操作のタイミングに加え、一手ごとの行動順と位置取りを考える体験に重点を置く。

\item[Contrast]
3Dの世界と2Dの影の世界を行き来し、光源や物体の操作で影を変えるパズル・プラットフォーム作品である\cite{contrast}。
本作品では、二匹と敵が共有するマップ上の明るさを変え、移動や敵の射撃条件へ影響させる。
光の操作を、予告された危険への対処と、次の攻撃の準備の両方に使う。

\item[Monster Sanctuary]
モンスターを育て、個別のスキルツリーとチームの相乗効果を利用して戦う作品である\cite{sanctuary}。
本作品では相棒を二匹に絞り、両者の明暗への適性と装備の効果を確認しながら育成する。
例えば、暗所で回避を高める装備は影を残す判断に、明所で攻撃を高める装備は照らす判断につながる設計とする。
技の成長分岐は授業での必須範囲に含めない。
\end{description}

\subsection{差別化と期待する利点}
\begin{enumerate}
\item 明暗操作により、攻撃・移動・敵の行動条件を変えられるため、直接攻撃以外にも状況を改善する手段がある。
\item 二匹の有利条件が異なるため、照らす場所と影を残す場所を考え、行動順を組み立てる必要がある。
\item 敵の予告を固定するため、操作前後の危険を比較し、自分の判断の結果を確認できる。回避は確率判定なので、予測では最大ダメージと回避率を分けて示す。
\end{enumerate}
以上は組合せと体験設計による差別化であり、既存作品全体に対する優越性や、個別要素が世界初であることを主張するものではない。

\section{詳細：具体的なソフトウェアの中身}
\subsection{プレイの流れと画面設計}
\begin{center}
\fbox{ホーム・育成}\ $\longrightarrow$\ \fbox{クエスト選択・条件確認}\\[2mm]
$\downarrow$\\[2mm]
\fbox{戦闘：行動順選択 $\rightarrow$ 二匹の行動 $\rightarrow$ 敵の行動}\\[2mm]
$\downarrow$\\[2mm]
\fbox{リザルト $\rightarrow$ 次の挑戦・育成・再挑戦}
\end{center}

\begin{tabularx}{\linewidth}{@{}p{25mm}Y@{}}
\toprule
画面 & 表示内容と操作 \\
\midrule
ホーム & 二匹の姿、レベル、経験値、進行状況、次の目的を表示する。クエスト選択と育成施設へ移動する。\\
クエスト選択 & エリアとクエストを選び、クリア条件・失敗条件を確認して出発する。\\
戦闘 & 目標、ラウンド、手番、HP、共通チャージ、明るさ、敵の予告を表示する。マスを選んで移動し、技と対象を選ぶ。\\
育成・装備 & 訓練の費用、所持素材、装備の効果と現在装備との差を示す。訓練、作成、付け替えを行う。\\
リザルト & 成功・失敗、獲得経験値・素材、成長と解放内容を表示する。次のクエスト、育成、ホーム、再挑戦へ進む。\\
\bottomrule
\end{tabularx}

\begin{figure}[htbp]
\centering
\fbox{\begin{minipage}{0.92\linewidth}
\small
上部：クエスト目標／ラウンド／手番／共通チャージ\\[2mm]
\begin{tabularx}{\linewidth}{|Y|p{43mm}|}
\hline
3Dマップ（内部は14$\times$8マス） & 二匹のHP・現在地の明るさ\\
移動範囲・経路・技の範囲 & 行動順／移動の確定・省略\\
敵の攻撃予告・移動予告 & 技の選択／対象／予測結果\\
\hline
\end{tabularx}
\end{minipage}}
\caption{PC横長画面の構成案。縦長画面では操作パネルをマップの下へ配置する。}
\end{figure}

移動可能範囲と経路を盤面に重ね、攻撃可能な対象と使用できない理由を表示する。
敵の攻撃予告には赤い斜線、移動予告には青い矢印を使い、色だけでなく形でも区別する。
光源や技の対象を選ぶと、実行後の明るさや被害予測を更新する。
影の猟犬は黒を基調にして暗所に溶け込む表現を目指す一方、選択表示や目・輪郭によって位置を確認できるようにする。
明るさの数値表示と演出軽減も用意する。

\subsection{戦闘・育成・進行の仕様}
\begin{tabularx}{\linewidth}{@{}p{27mm}Y@{}}
\toprule
項目 & 内容 \\
\midrule
マップ・移動 & 横14、縦8マス。上下左右に移動し、生存キャラクターや遮蔽物を通り抜けない。\\
ラウンド & 開始時に光から／影からを選ぶ。各犬が一回ずつ移動後に技を一つ使い、その後に敵が行動する。移動と技はそれぞれ省略できる。\\
光の猟犬 & 通常攻撃、灯火の移動、直線状の光の技、防御を行う。明所で攻撃や移動が有利になる。\\
影の猟犬 & 通常攻撃、遮蔽物の設置、直線状の影の技、防御を行う。暗所で移動・攻撃・回避が有利になる。\\
敵の予告 & ラウンド開始時に攻撃マス・移動先などを固定する。敵手番で狙い直さず、遮蔽物や明暗の変化で実行できなくなれば中止する。\\
共通チャージ & 最大3、開始時1、敵手番終了後に1回復。光・影の直線技と遮蔽物設置は2消費する。直線技は明暗を変え、直接ダメージを与えない。\\
クエスト & 討伐、秘宝の回収、迷子の救助、ボス討伐の4種。回収・救助は対象を拾った猟犬が帰還して達成する。\\
育成・装備 & 二匹が個別のHP・経験値・レベルを持つ。素材による攻撃・体力訓練と装備の作成・付け替えを行う。出発時の能力を戦闘に反映する。\\
終了・再挑戦 & 一匹が倒れても続行し、二匹とも倒れると失敗。成功時は報酬を一度だけ付与する。失敗後の再挑戦ではHPを全回復する。\\
保存 & localStorageへ育成、素材、装備、達成記録、設定を保存する。進行中の戦闘は保存しない。\\
\bottomrule
\end{tabularx}

\subsection{明るさのルール}
明るさは0から4で管理する。光源との距離と遮蔽物から算出し、灯火の移動、遮蔽物の設置、技の使用や効果終了に合わせて再計算する。
複数の光源がある場合は各光源から得られる明るさの最大値を使い、最後に技の効果を適用する。
以下は現行試作の調整値であり、試遊結果を基に調整する。

\begin{center}
\small
\begin{tabular}{@{}clrrc@{}}
\toprule
明るさ & 状態 & 光の移動予算 & 影の移動予算 & 影の進入コスト\\
\midrule
0 & 濃い影 & 1 & 5 & 1\\
1 & 暗がり & 2 & 4 & 2\\
2 & 薄明 & 3 & 3 & 3\\
3 & 明るい & 3 & 2 & 3\\
4 & 強い光 & 4 & 0 & 進入不可\\
\bottomrule
\end{tabular}
\end{center}

移動予算は移動開始時のマスの明るさで決まり、経路の各マスへの進入コストの合計が予算以内なら移動できる。
光の進入コストは1である。
影の通常攻撃は、明るさ4の敵には使用できない。
影の現在地が4になった場合は有効な暗いマスへ退避し、行動を一回失う。
退避先が存在しない場合の扱いは未確定事項として、10月22日の企画書提出までに方針を整理し、11月6日の試作確認までに確定して検証項目に加える。

\subsection{使用する言語・ライブラリと選定理由}
\begin{tabularx}{\linewidth}{@{}p{31mm}Y@{}}
\toprule
技術 & 用途と採用理由 \\
\midrule
TypeScript／JavaScript & 状態、能力、装備、戦闘などの処理を実装する。新規ルールはTypeScriptで記述し、状態や引数の型を確認する。既存JavaScriptは動作を維持しながら段階的に移行する。\\
HTML／CSS & ボタン、説明、HP、ホーム、リザルトを実装する。画面幅に応じて配置を変え、マウスとタッチの両方で操作できるUIを作る。\\
Three.js／WebGL & カメラ、犬・地形・装備の3D描画、照明を扱う。導入済みの描画基盤を利用し、2Dのマス座標と3Dの見た目を分離する。\\
Vite／npm & 開発サーバーと配布用ビルドに用いる。package-lock.jsonを共有し、npm ciで依存関係をそろえる。\\
localStorage & サーバーを設けずブラウザ内へ進行を保存する。同じURLの同じブラウザでの利用を前提とし、端末間の同期は扱わない。\\
Web Audio API & 操作・移動・攻撃などの効果音を合成する。最初のユーザー操作後に音声を有効にし、音量・ミュートを設定できるようにする。\\
Vitest／Playwright & 前者でゲームルール、後者でブラウザ上の画面遷移・入力・保存を検証する。実機でのタッチ操作・音声確認も別途行う。\\
\bottomrule
\end{tabularx}

配布はビルドしたファイルをHTTPS対応の静的ホスティングへ置く方式を予定する。公開先は未定である。
対戦用サーバー、データベース、クラウド保存は完成範囲に含めない。

\subsection{ソフトウェアの構成と現在地}
ゲームルール、UI、描画、音声、保存の責務を分ける。
UIが操作を受け取り、ゲームルールが状態を更新し、その状態を描画・表示へ渡す構成を目指す。
明るさによる能力値は共通の定義から取得し、敵の予測表示と実際の命中も同じ判定処理を使うことで、表示と結果のずれを防ぐ。

現在のコードにはVite、Three.js、TypeScriptの一部ルール、装備、クエスト進行、行動順・敵予告、テストが存在する。
クエストは3エリア・計15件の定義があり、4種の達成条件を利用する。
一方、既存JavaScriptに統合処理が残り、全体のTypeScript移行は未完了である。
また、コードやテストの存在だけで、実機での操作性や難易度の検証が完了したとは扱わない。

\section{制作計画}
\subsection{完成範囲と優先順位}
既存の試作を基に、二匹の操作、明暗操作、敵の固定予告、4種のクエスト、基本育成、装備、ホーム、リザルト、保存を一連の流れとして完成させる。
まず各種一件以上の代表クエストで開始から報酬・保存までを通し、その後、既存15クエストの配置・解放・難易度を確認する。
固定能力・報酬なしの戦術演習は、実際の育成データを変えずに連携を練習する導線として整える。

高さによる射線・移動、技の成長分岐、対人戦、クラウド保存、外部リグ付きモデルの導入は今回の必須範囲から外す。
期限内の完成を優先し、遅れた場合は追加コンテンツとモデルの細部改善を後回しにして、基本の遊びと保存の安定性を確保する。

\subsection{役割分担（仮案）}
担当の希望・経験は未確認のため、以下は作業の境界を明確にするための仮案とする。
全員が機能実装と担当部分の検証コードを受け持ち、素材作成のみを担当するメンバーは設けない。
本企画は村松のアイデアを中心に始まったため、村松が企画意図の共有と仕様調整も担う案とする。
追加の実装や統合作業を村松が多めに担当する可能性はあるが、作業量は週ごとの進捗を見て調整し、他の三名も担当機能の実装を受け持つ。

\begin{tabularx}{\linewidth}{@{}p{13mm}Y Y@{}}
\toprule
担当 & 主なプログラミング作業 & 成果物と検証内容 \\
\midrule
木澤 & 戦闘・手番管理。行動順、移動後の技、攻撃、敵の固定予告、勝敗判定の実装・整理。 & 一ラウンドの状態遷移。各犬一回の行動、予告と命中の一致、戦闘不能時のテスト。\\
村松 & マップ・明暗・経路。光源と遮蔽物の判定、移動可能範囲、明暗操作後の予測計算、クエスト配置。 & 明るさと移動結果を返す処理。進入不可、遮蔽、影の退避、経路コストのテスト。\\
中野 & 3D描画・UI。Three.jsの描画更新とマス選択、犬の動作、予告・範囲表示、画面配置と入力処理。 & PC・縦画面の操作UI。マス選択、画面遷移、演出中の重複入力をブラウザで検証。\\
森井 & 育成・装備・進行・保存。経験値、素材消費、装備効果、報酬、解放条件、保存読込の実装・整理。 & クリアから育成・次の出発までの処理。報酬の一回付与、能力反映、旧保存データのテスト。\\
\bottomrule
\end{tabularx}

木澤と村松は盤面状態・予測計算の受け渡し、中野と各担当は表示に渡す状態と操作の形式、森井と木澤は出発時の能力・終了時の報酬の境界を合わせる。
共通の状態形式と関数の入出力は10月23日の試作着手時に共有し、既存の統合ファイルを複数人が同時に大きく書き換えないよう、担当機能から順に分離する。
週1回の進捗共有と統合確認を行い、各自が実装・検証結果・次の作業を共有する。
試遊と難易度調整は全員で行う。

\subsection{制作スケジュール}
授業日程に合わせ、企画書は2026年10月22日（木）20時までに提出し、11月13日の中間発表で企画を確定する。
制作物の完成目標は12月18日とし、2027年1月15日に体験形式、1月22日にプレゼン形式の最終発表を行う。
以下の授業・提出日程に対し、各期間の作業と到達目標を計画する。

\begin{tabularx}{\linewidth}{@{}p{25mm}Y Y@{}}
\toprule
期間・授業日 & 主な作業 & 到達目標 \\
\midrule
10月16日～22日 & 全員：企画書を作成し、完成範囲、担当案、技術構成を共有する。村松を中心に企画意図と未確定仕様を整理する。 & 10月22日（木）20時までに企画書を提出する。\\
10月23日～11月6日 & 木澤：手番・予告。村松：明暗・移動。中野：入力・範囲表示。森井：報酬・進行。既存試作を確認して担当機能を統合する。 & 討伐一件を開始からリザルトまで通す。11月6日に試作確認と企画書の再提出へ対応する。\\
11月13日：中間発表 & 全員：操作デモと進捗を発表する。実装済み・未完了を整理し、完成範囲と分担を見直す。 & 企画書を確定し、明暗操作と二匹の連携を示せる状態にする。\\
11月20日：制作 & 木澤・村松：4種の達成条件と予告への対処。中野：目標・運び手・予測表示。森井：各条件と報酬の接続。 & 各種一件以上で、開始から達成・報酬まで通して遊べる。\\
11月27日：制作 & 森井：育成・装備・保存。木澤：出発時能力との接続。村松：装備込みの移動とクエスト配置。中野：育成・装備UI。 & 育成を次の出発へ反映し、再読込後も保持する。既存15件の条件と解放を点検する。\\
12月4日：制作 & 中野：犬の動作・明暗表現・画面配置。全員：演習、PC・スマホ操作、音声、連打と再挑戦を確認する。 & 主要機能をそろえ、以後は機能追加を止める。見た目と判定の不一致を修正する。\\
12月11日：制作 & 全員：チーム外の人の試遊を観察し、案内・難易度・二匹の使い分けを調整する。自動検証と実機確認を行う。 & 理解できなかった点と修正を記録し、進行不能や保存の不具合を解消する。\\
12月18日：完成 & 全員：回帰検証、公開先での動作確認、既知の制約の整理、デモ手順と予備動画の作成を行う。 & 体験可能な配布版を完成させ、発表に使う版を保存する。\\
1月13日～15日 & 全員：休み明けの起動・保存・公開URLを再確認し、体験時の案内を準備する。 & 1月15日の最終発表で制作物を体験できるようにする。\\
1月16日～22日 & 全員：作品の目的、実装した仕組み、分担、検証結果を発表資料にまとめる。 & 1月22日にプレゼン形式の最終発表を行う。\\
\bottomrule
\end{tabularx}

10月30日は文化祭準備のため休講であり、授業内の統合作業は予定しない。
12月23日から1月12日までの休み中の作業を完成の前提にせず、12月18日までに制作物を完成させる。
10月23日以降は毎週、担当部分を統合して代表クエストを通す。
遅れが生じた場合は中間発表および12月4日の時点で優先順位を見直し、12月11日以降の試遊・不具合修正の時間を確保する。

\subsection{評価・完成確認の方法}
初めて遊ぶ人に、目標確認、灯火の移動、二匹の連携、クリア後の育成までを試してもらう。
操作が止まった箇所、誤解したルール、説明を必要とした場面を記録する。
その後、光を動かす意味、二匹を使い分ける理由、敵予告への対処、帰還条件を本人の言葉で説明できるか確認し、UIと案内の修正に使う。
試遊だけで面白さや愛着の向上を証明したとは扱わない。

自動検証では、明るさと能力の対応、行動順、予告と命中、回収者の帰還、報酬の一回付与、素材消費、保存後の復元を確認する。
ビルド時の型検査、Vitestのルール検証、Playwrightの画面操作を実施し、PC・スマホ実機ではタッチの選択位置、画面の欠け、音量と効果音を別に確認する。
外部素材を追加する場合は出典・ライセンスを記録し、未解決の制約は配布版の説明に残す。

\begin{thebibliography}{9}
\bibitem{puppeteer}
Sarepta studio, Shadow Puppeteer, Steam作品紹介ページ，
\url{https://store.steampowered.com/app/316480/Shadow_Puppeteer/}
（参照日：2026年10月16日）．
\bibitem{contrast}
Compulsion Games, Contrast, Steam作品紹介ページ，
\url{https://store.steampowered.com/app/224460/Contrast/}
（参照日：2026年10月16日）．
\bibitem{sanctuary}
Moi Rai Games, Monster Sanctuary, Steam作品紹介ページ，
\url{https://store.steampowered.com/app/814370/Monster_Sanctuary/}
（参照日：2026年10月16日）．
\end{thebibliography}
\end{document}
